% !TEX program = xelatex
\documentclass[UTF8,a4paper,11pt,fontset=fandol]{ctexart}

\usepackage[margin=2.2cm]{geometry}
\usepackage{xcolor}
\usepackage{listings}
\usepackage{amsmath}
\usepackage{hyperref}
\usepackage{enumitem}

\hypersetup{
    colorlinks=true,
    linkcolor=black,
    urlcolor=blue
}

\lstdefinestyle{cppstyle}{
    language=C++,
    basicstyle=\ttfamily\small,
    keywordstyle=\color{blue!70!black},
    commentstyle=\color{green!40!black},
    stringstyle=\color{orange!70!black},
    numbers=left,
    numberstyle=\tiny\color{gray},
    frame=single,
    breaklines=true,
    tabsize=4,
    showstringspaces=false,
    columns=fullflexible,
    keepspaces=true
}

\title{Lab 3 FEM 软体与布料仿真实现报告}
\author{}
\date{}

\begin{document}
\maketitle

\section{软体模拟部分说明}

软体模拟的核心代码位于 \texttt{src/VCX/Labs/3-FEM/FEMSystem.h/.cpp} 和
\texttt{CaseSoftBody.h/.cpp}。系统使用规则长方体网格生成四面体网格，默认尺寸为
$8\mathrm{m}\times 2\mathrm{m}\times 2\mathrm{m}$，网格分辨率为 $8\times2\times2$。
每个立方体单元被拆成 6 个四面体；左端 $i=0$ 的顶点固定，其余顶点作为自由质点参与显式积分。

每个四面体在初始化时保存四个顶点索引、Rest Shape 矩阵逆 $\mathbf{D}_m^{-1}$ 和静止体积 $V$。
其中
\[
\mathbf{D}_m =
\begin{bmatrix}
\mathbf{X}_1-\mathbf{X}_0 &
\mathbf{X}_2-\mathbf{X}_0 &
\mathbf{X}_3-\mathbf{X}_0
\end{bmatrix},
\qquad
V=\frac{|\det(\mathbf{D}_m)|}{6}.
\]
质量按体积和密度计算，再平均分配到四面体的四个顶点。固定点质量设为无穷大，并在每步积分时强制回到 Rest Position。

\begin{lstlisting}[style=cppstyle,caption={四面体初始化时记录 Rest Shape 信息}]
void FEMSystem::AddTet(std::size_t i0, std::size_t i1,
                       std::size_t i2, std::size_t i3) {
    glm::mat3 Dm = makeColumnMat(
        RestPositions[i1] - RestPositions[i0],
        RestPositions[i2] - RestPositions[i0],
        RestPositions[i3] - RestPositions[i0]);
    float det = glm::determinant(Dm);
    if (det < 0.f) {
        std::swap(i1, i2);
        Dm = makeColumnMat(
            RestPositions[i1] - RestPositions[i0],
            RestPositions[i2] - RestPositions[i0],
            RestPositions[i3] - RestPositions[i0]);
        det = glm::determinant(Dm);
    }

    Tets.push_back(Tet {
        .Idx        = { i0, i1, i2, i3 },
        .InvRest    = glm::inverse(Dm),
        .RestVolume = std::abs(det) / 6.f,
    });
}
\end{lstlisting}

每一小步仿真先清空力，再加入重力、速度阻尼和用户外力。对每个四面体，当前形状矩阵为
\[
\mathbf{D}_s =
\begin{bmatrix}
\mathbf{x}_1-\mathbf{x}_0 &
\mathbf{x}_2-\mathbf{x}_0 &
\mathbf{x}_3-\mathbf{x}_0
\end{bmatrix},
\qquad
\mathbf{F}=\mathbf{D}_s\mathbf{D}_m^{-1}.
\]
基础模型使用 StVK 能量。记
\[
\mathbf{G}=\frac{1}{2}(\mathbf{F}^{T}\mathbf{F}-\mathbf{I}),\qquad
\mathbf{S}=2\mu \mathbf{G}+\lambda \operatorname{tr}(\mathbf{G})\mathbf{I},
\]
则 First Piola--Kirchhoff stress 为
\[
\mathbf{P}=\mathbf{F}\mathbf{S}.
\]
离散力使用
\[
\mathbf{H}=-V\mathbf{P}\mathbf{D}_m^{-T}.
\]
\texttt{H} 的三列分别加到第 1、2、3 个顶点，第 0 个顶点受到负的列和，从而保证单元内部力守恒。

\begin{lstlisting}[style=cppstyle,caption={由 deformation gradient 计算四面体内力}]
glm::mat3 const Ds = makeColumnMat(
    Positions[i1] - Positions[i0],
    Positions[i2] - Positions[i0],
    Positions[i3] - Positions[i0]);
glm::mat3 const F = Ds * tet.InvRest;
glm::mat3 const P = firstPiola(F, Model, lambda, mu);
glm::mat3 const H = -tet.RestVolume * P * glm::transpose(tet.InvRest);

Forces[i1] += col(H, 0);
Forces[i2] += col(H, 1);
Forces[i3] += col(H, 2);
Forces[i0] -= col(H, 0) + col(H, 1) + col(H, 2);
\end{lstlisting}

时间积分采用显式欧拉。自由顶点根据 $\mathbf{a}_i=\mathbf{f}_i/m_i$ 更新速度和位置，并对速度做上限裁剪以避免显式积分在极端参数下数值爆炸。默认物理参数采用作业文档推荐值：杨氏模量 $20000$，泊松比 $0.2$，密度 $400$，重力 $0.05$，阻尼 $1.2$，时间步长 $0.001$。场景默认停止，需要点击 \texttt{Start Simulation} 才开始仿真。

\begin{lstlisting}[style=cppstyle,caption={显式欧拉积分与速度上限}]
glm::vec3 const accel = Forces[i] / Masses[i];
Velocities[i] += accel * dt;
float const speed = glm::length(Velocities[i]);
if (speed > MaxVelocity) {
    Velocities[i] *= MaxVelocity / speed;
}
Positions[i] += Velocities[i] * dt;
\end{lstlisting}

为了让仿真速度不直接依赖帧率，软体部分在渲染时使用 \texttt{Engine::GetDeltaTime()} 做时间累积。每帧按照固定小步长 \texttt{\_timeStep} 消耗累积时间，同时通过 \texttt{Max Steps / Frame} 限制单帧最多子步数。

\begin{lstlisting}[style=cppstyle,caption={软体模拟的固定时间步推进}]
float const frameDt = std::min(Engine::GetDeltaTime(), 0.05f);
_timeAccumulator += frameDt * _simulationSpeed;

int steps = 0;
int const maxSteps = std::max(1, _maxStepsPerFrame);
while (_timeAccumulator >= _timeStep && steps < maxSteps) {
    _system.Step(_timeStep, _selectedParticle, frameControlForce);
    _timeAccumulator -= _timeStep;
    ++steps;
}
\end{lstlisting}

\section{B1 说明：多种超弹性模型}

Bonus 1 在同一个软体系统中加入材料模型选择，而不是同时显示多个软体。代码中定义了
\texttt{MaterialModel} 枚举，并在 UI 中通过 \texttt{Material Model} 下拉框切换
\texttt{StVK}、\texttt{Neo-Hookean} 和 \texttt{Corotated}。切换后同一个四面体网格会使用新的应力模型计算内力，因此可以在相同几何和交互条件下观察材料行为差异。

\begin{lstlisting}[style=cppstyle,caption={材料模型枚举与 UI 选择}]
enum class MaterialModel : int {
    StVK = 0,
    NeoHookean = 1,
    Corotated = 2,
};

char const * names[] = { "StVK", "Neo-Hookean", "Corotated" };
if (ImGui::Combo("Material Model", &_activeModel, names, 3)) {
    _system.Model = static_cast<MaterialModel>(_activeModel);
}
\end{lstlisting}

三种模型共用同一套四面体离散和力汇总流程，区别只在 \texttt{firstPiola} 函数。StVK 模型使用上一节所述
$\mathbf{P}=\mathbf{F}\mathbf{S}$。Neo-Hookean 模型使用
\[
\mathbf{P}
= \mu(\mathbf{F}-\mathbf{F}^{-T})
  + \lambda \log(J)\mathbf{F}^{-T},
\qquad J=\det(\mathbf{F}).
\]
实现中对 $J$ 做了最小值裁剪，避免单元翻转或过度压缩时 $\log(J)$ 和逆矩阵产生非法值。

Corotated 模型先对 $\mathbf{F}$ 做极分解，取旋转部分 $\mathbf{R}$，再使用
\[
\mathbf{P}
= 2\mu(\mathbf{F}-\mathbf{R})
  + \lambda (J-1)J\mathbf{F}^{-T}.
\]
极分解通过 Eigen 的 SVD 实现，并在 $\det(\mathbf{R})<0$ 时修正一列奇异向量，保证得到合法旋转矩阵。

\begin{lstlisting}[style=cppstyle,caption={三种材料模型的 First Piola stress}]
glm::mat3 firstPiola(glm::mat3 const & F,
                     MaterialModel model,
                     float lambda,
                     float mu) {
    if (model == MaterialModel::NeoHookean) {
        float const J = std::max(glm::determinant(F), 1e-4f);
        glm::mat3 const FinvT = glm::transpose(glm::inverse(F));
        return mu * (F - FinvT) + lambda * std::log(J) * FinvT;
    }

    if (model == MaterialModel::Corotated) {
        glm::mat3 const R = polarRotation(F);
        float const J = glm::determinant(F);
        return 2.f * mu * (F - R)
             + lambda * (J - 1.f) * J * glm::transpose(glm::inverse(F));
    }

    glm::mat3 const C = glm::transpose(F) * F;
    glm::mat3 const G = 0.5f * (C - glm::mat3(1.f));
    glm::mat3 const S = 2.f * mu * G
                      + lambda * traceMat(G) * glm::mat3(1.f);
    return F * S;
}
\end{lstlisting}

\section{B2 说明：布料模拟}

Bonus 2 的布料代码位于 \texttt{ClothSystem.h/.cpp} 和 \texttt{CaseCloth.h/.cpp}。布料被建模为嵌入三维空间的二维三角形网格，默认分辨率为 $24\times24$，尺寸为 $2\mathrm{m}\times2\mathrm{m}$。网格初始位于 $xz$ 平面上方，左侧两个角固定，其余顶点自由运动。

和三维软体不同，布料单个三角形的 Rest Shape 只需要二维坐标描述。初始化时为每个顶点保存三维位置 \texttt{RestPositions} 和二维参数坐标 \texttt{RestUV}。每个三角形保存
\[
\mathbf{D}_m =
\begin{bmatrix}
\mathbf{u}_1-\mathbf{u}_0 &
\mathbf{u}_2-\mathbf{u}_0
\end{bmatrix},
\qquad
A=\frac{|\det(\mathbf{D}_m)|}{2}.
\]
质量根据面密度和三角形面积分配到三个顶点，默认参数使用文档建议值：杨氏模量 $50$，泊松比 $0.3$，面密度 $0.5$，重力 $9.8$，时间步长 $0.001$。为了让默认演示更稳定，速度阻尼设为 $0.6$，默认每帧只推进 1 个子步；需要更快运动时可以手动增加 \texttt{Steps / Frame}。

布料的当前形状矩阵为 $3\times2$：
\[
\mathbf{D}_s =
\begin{bmatrix}
\mathbf{x}_1-\mathbf{x}_0 &
\mathbf{x}_2-\mathbf{x}_0
\end{bmatrix},
\qquad
\mathbf{F}=\mathbf{D}_s\mathbf{D}_m^{-1}.
\]
随后在二维参数域中计算
\[
\mathbf{C}=\mathbf{F}^{T}\mathbf{F},\qquad
\mathbf{G}=\frac{1}{2}(\mathbf{C}-\mathbf{I}_2),
\]
\[
\mathbf{S}=2\mu\mathbf{G}+\lambda\operatorname{tr}(\mathbf{G})\mathbf{I}_2,
\qquad
\mathbf{P}=\mathbf{F}\mathbf{S}.
\]
离散力矩阵为
\[
\mathbf{H}=-A\mathbf{P}\mathbf{D}_m^{-T}.
\]
\texttt{H} 是 $3\times2$ 矩阵，其两列分别给三角形第 1、2 个顶点，第 0 个顶点受到负列和。

\begin{lstlisting}[style=cppstyle,caption={布料三角形的 3x2 FEM 内力}]
glm::mat2x3 const Ds(
    Positions[i1] - Positions[i0],
    Positions[i2] - Positions[i0]);
glm::mat2x3 const F = Ds * tri.InvRest;
glm::mat2 const C = glm::transpose(F) * F;
glm::mat2 const G = 0.5f * (C - glm::mat2(1.f));
glm::mat2 const S = 2.f * mu * G
                  + lambda * traceMat(G) * glm::mat2(1.f);
glm::mat2x3 const P = F * S;
glm::mat2x3 const H = -tri.RestArea * P
                    * glm::transpose(tri.InvRest);

glm::vec3 const f1 = mat2x3Col0(H);
glm::vec3 const f2 = mat2x3Col1(H);
Forces[i1] += f1;
Forces[i2] += f2;
Forces[i0] -= f1 + f2;
\end{lstlisting}

除了面内 FEM 弹性力，布料还加入了简单的弯曲弹簧。实现上对横向或纵向相隔两个网格点的顶点建立 \texttt{BendEdge}，保存静止长度。每步对这些边施加 Hooke 型长度恢复力，用于减小布料过度折叠和数值抖动。

\begin{lstlisting}[style=cppstyle,caption={布料弯曲弹簧}]
glm::vec3 const x = Positions[edge.I1] - Positions[edge.I0];
float const len = glm::length(x);
if (len > 1e-7f) {
    glm::vec3 const dir = x / len;
    glm::vec3 const f = BendingStiffness
                      * (len - edge.RestLength) * dir;
    if (!Fixed[edge.I0]) Forces[edge.I0] += f;
    if (!Fixed[edge.I1]) Forces[edge.I1] -= f;
}
\end{lstlisting}

\section{交互方式指南}

\begin{itemize}[leftmargin=*]
    \item 构建运行：在仓库根目录执行 \texttt{xmake run lab3}。
    \item 软体材料：软体场景中通过 \texttt{Material Model} 下拉框切换 \texttt{StVK}、\texttt{Neo-Hookean}、\texttt{Corotated}。
    \item 选择控制点：\texttt{Control Particle} 滑条可以选择任意顶点；软体的 \texttt{Select Free End} 会自动选择自由端顶点，布料的 \texttt{Select Free Corner} 会自动选择自由角点。
    \item 键盘施力：画面获得焦点后，\texttt{J/L} 沿 $-X/+X$ 方向施力，\texttt{U/O} 沿 $+Y/-Y$ 方向施力，\texttt{I/K} 沿 $-Z/+Z$ 方向施力。
    \item 鼠标拖拽：按住 \texttt{Z + 鼠标左键} 并拖动画面，可以在当前相机视平面内对选中顶点施加外力。拖拽强度由 \texttt{Mouse Force Scale} 控制。
    \item 按钮施力：UI 中的 \texttt{+X/-X/+Y/-Y/+Z/-Z} 按钮会对当前选中顶点施加一次方向力，大小由 \texttt{Control Force} 控制。
    \item 可视化：黄色大点表示当前选中的控制顶点；红色点表示固定顶点；\texttt{Surface}、\texttt{Wireframe}、\texttt{Particles}、\texttt{Fixed Points} 可分别控制表面、线框、粒子和固定点显示。
\end{itemize}

\end{document}
